Un estudi ha revelat que el contingut de cesi 137 en vins de diferents anyades es relaciona amb diferents episodis nuclears, especialment amb les explosions termonuclears que van ocórrer durant la guerra freda: el Cs-137 emès en aquests episodis es va incorporar als vins, i un cop a l'ampolla no hi va haver cap entrada ni sortida d'aquest isòtop radioactiu. Aquest isòtop es desintegra emetent radiacions beta i gamma que poden ser detectades i mesurades sense haver d'obrir les ampolles.

La gràfica mostra la radioactivitat que tenien els vins de cada any en el moment de la vinificació. La línia discontínua mostra l'activitat dels vins de la collita del 1957 a cada moment.

\figura[0.55]{fig1.png}

\begin{apartat}{1,25}
Segons la gràfica els vins embotellats el 1957 tenien inicialment una activitat de 2\,500~mBq (cada litre), que s'ha reduït fins a 600~mBq/l l'any 2020. Amb aquestes dades, calculeu el període de semidesintegració del Cs-137 i dibuixeu, sobre la mateixa gràfica de l'enunciat, la corba de decaïment per als vins embotellats el 1972 (any de les darreres proves nuclears franceses no subterrànies).
\end{apartat}

\begin{apartat}{1,25}
El Cs-137 es produeix a partir de processos com el que es representa a aquesta equació nuclear: ${}^{235}_{92}\mathrm{U} + {}^{1}_{0}? \rightarrow {}^{137}_{55}\mathrm{Cs} + {}^{x}_{y}\mathrm{Rb} + 3\,{}^{1}_{0}?$

Completeu l'equació indicant els valors de $x$ i $y$, i la naturalesa de la partícula indicada amb un interrogant. Digueu també de quina classe de reacció nuclear es tracta.
\end{apartat}

\blocetiquetat{Font:}{Michael S. \textsc{Pravikoff} i Philippe \textsc{Hubert}, \textit{Dating of wines with cesium-137: Fukushima's imprint}, <arXiv:1807.04340>.}
